% SPDX-FileCopyrightText: Copyright (c) 2016-2026 Objectionary.com
% SPDX-License-Identifier: MIT

\section{Operators}\label{sec:operators}

In this section we introduce a family of semantic operators
  defined on \(\phi\)-expressions.
These operators do not extend the syntax of the calculus.
Instead, they describe systematic transformations of expressions that reveal,
  modify, or project their semantic content.
Each operator is defined in terms of the reduction semantics introduced earlier
  and preserves the core meaning of expressions
  while possibly changing their form or role.

\subsection{Contextualization}\label{sec:contextualization}
\makeatletter
\protected@edef\@currentlabelname{Ctx}%
\makeatother
\label{r:contextualize}

\begin{definition}[Universe]
A \textbf{universe} is an expression that serves as the global context of
  another expression, in which the $Q$ locator refers to the universe.
\end{definition}

\begin{definition}[Absolute]
An expression is \textbf{absolute},
  denoted by \(k\) and ranging over \(\mathcal{K} \subseteq \mathcal{E}\),
  if it is either
  \begin{inparaenum}[a)]
  \item $Q$,
  \item formation,
  \item dispatch with an absolute subject,
  or
  \item application with absolute subject and argument.
  \end{inparaenum}
\end{definition}

\begin{definition}[Contextualization]
The \emph{contextualization} function
  \(\mathcal{E} \times \mathcal{K} \to \mathcal{E}\),
  denoted as \(\phinoContextualize{e_1}{k}{e_2}\)
  and specified by the rules of \cref{fig:contextualization},
  turns an expression \(e_1\) into \(e_2\)
  by replacing with the absolute \(k\)
  every \(\phiTerminal{\xi}\) that denotes the context of \(e_1\)
  rather than that of a nested formation.
\end{definition}

\begin{figure*}
\begin{mdframed}
\centering
\iexec[maybe]{phino explain --contextualize > _tex/rules-contextualize.tex}
\input{_tex/rules-contextualize.tex}
\end{mdframed}
\capt{Contextualization by induction.}
\label{fig:contextualization}
\end{figure*}

\begin{theorem}\label{thm:contextualization-total}
The contextualization function is total, meaning that
  every expression can be contextualized and its derivation terminates.
\end{theorem}

\subsection{Normalization}\label{sec:normalization}

An expression that may be rewritten by the \emph{rules} (or \emph{reductions})
  listed in \cref{fig:reduction} is a \emph{reducible} expression.
The notation \(e_1 \trans e_2\), optionally followed by a condition,
  denotes a reduction of \(e_1\) to \(e_2\), if the condition holds.

A specific reduction may be denoted, for example,
  as \(\trans_{\nameref{r:dot}}\),
  or just \(\trans\) when no specific reduction is meant.
The notation \(e_1 \strans e_2\) denotes a reflexive transitive
  closure of all reductions, so that there is a possibly empty finite
  sequence of reductions between \(e_1\) and \(e_2\).

\begin{definition}[Normal Form]
An expression that has no more possible applications of reductions
  is \emph{irreducible} or a \emph{normal form}, denoted as \(\nf{}_i\)
  ranging over \(\mathcal{N} \subseteq \mathcal{E}\).
\end{definition}

Thus, \(\nf\) is a normal form of \(e\) if \(e \strans \nf\) and
  there is no expression \(e_1\) such that \(\nf \trans e_1\).

\begin{figure*}
\begin{mdframed}
\renewcommand{\arraystretch}{1.2}
\iexec[maybe]{phino explain --normalize > _tex/rules.tex}
\inputrules{_tex/rules.tex}
\end{mdframed}
\capt{Reduction rules, generated by phino \iexec{phino --version}\unskip.}
\label{fig:reduction}
\end{figure*}

Not every expression has a normal form.
For some expressions normalization may never terminate, for example:

\iexec[maybe]{echo '[[ x -> $.y, y -> $.x ]].x' | phino rewrite --normalize --max-depth=5 --max-cycles=1 --nonumber --sequence --compress --output=latex --flat --sweet}

\Cref{app:normalization-examples} demonstrates how normalization
  works through examples.

\begin{theorem}\label{thm:normalization-confluent}
The reduction relation is \emph{confluent}:
  if \(e \strans e_1\) and \(e \strans e_2\),
  then there exists an expression \(e_3\) such that
  \(e_1 \strans e_3\) and \(e_2 \strans e_3\).
\end{theorem}

Confluence guarantees that the normal form of an expression is unique,
  irrespective of the order in which the rules are applied.

\subsection{Equivalence}\label{sec:equivalence}

\begin{definition}[Equivalence]
Two expressions are said to be syntactically equivalent or just
  \textbf{equivalent} (denoted by \(\equiv\)) if their normal forms are
  syntactically identical.
\end{definition}

\subsection{Evaluation}\label{sec:evaluation}
\makeatletter
\protected@edef\@currentlabelname{Evt}%
\makeatother
\label{r:evaluate}

\begin{definition}[Evaluation]
The \emph{evaluation} function
  \(\mathcal{B} \times \mathcal{E} \times \mathcal{S}
    \rightharpoonup \mathcal{E} \times \mathcal{S} \),
  denoted as \(\phinoEvaluate{b}{e}{s_1}{n}{s_2}\),
  maps an atom \(b\) to \(n\),
  calling by value the function attached to its $L$-asset
  in the universe \(e\),
  possibly shifting the state \(s_1\) to \(s_2\):
\begin{phiquation*}
\dfrac{F \lparen [[ B_1, L>F, B_2 ]] {,}\; e {,}\; s_1 \rparen \rightarrow \langle n {,}\; s_2 \rangle}{\phinoEvaluate{[[ B_1, L>F, B_2 ]]}{e}{s_1}{n}{s_2}}
\end{phiquation*}
\end{definition}

\subsection{Morphing}\label{sec:morphing}

\begin{definition}[Morphing]
The \emph{morphing} function
  \(\mathcal{N} \times \mathcal{E} \times \mathcal{S}
    \rightharpoonup ( \mathcal{B} \cup \{ T \} ) \times \mathcal{S} \),
  denoted as \(\phinoMorph{n_1}{e}{s_1}{n_2}{s_2}\)
  and specified by the rules of \cref{fig:morphing},
  reduces a normal form \(n_1\) to the formation it denotes,
  or to $T$ if it denotes none,
  by evaluating the atoms it reaches
  and resolving the $Q$ locator against the universe \(e\),
  possibly shifting the state \(s_1\) to \(s_2\).
Morphing is a partial function, since it normalizes the expressions it
  reaches (\cref{sec:normalization}) and normalization may not terminate.
\end{definition}

\begin{figure*}
\begin{mdframed}
\centering
\iexec[maybe]{phino explain --morph > _tex/rules-morph.tex}
\input{_tex/rules-morph.tex}
\end{mdframed}
\capt{Morphing rules, generated by phino \iexec{phino --version}\unskip.}
\label{fig:morphing}
\end{figure*}

\begin{theorem}\label{thm:morphing-confluent}
The morphing function is confluent, meaning that
  every normal form can be morphed to at most one formation or $T$.
\end{theorem}

\subsection{Dataization}\label{sec:dataization}

\begin{definition}[Dataization]
The \emph{dataization} function
  \(\mathcal{N} \times \mathcal{E} \times \mathcal{S}
    \rightharpoonup \mathcal{D} \times \mathcal{S}\),
  denoted as \(\phinoDataize{n}{e}{s_1}{\delta}{s_2}\)
  and specified by the rules of \cref{fig:dataization},
  extracts the datum \(\delta\) that a normal form \(n\) produces
  by morphing \(n\) to a formation in the universe \(e\)
  and reading the data attached to its $D$-asset,
  possibly shifting the state \(s_1\) to \(s_2\).
\end{definition}

\begin{theorem}[Progress]\label{thm:dataization-progress}
Dataization is defined for every normal form except the terminator $T$,
  meaning that for any \(n \neq T\), \(e\), and \(s_1\) at least one rule of
  \cref{fig:dataization} applies, so a derivation never gets stuck at a term
  that ought to carry data.
The terminator $T$ matches no rule: a derivation that reaches it halts without
  a datum, which is exactly the error that $T$ signals.
\end{theorem}

\begin{figure*}
\begin{mdframed}
\centering
\iexec[maybe]{phino explain --dataize > _tex/rules-dataize.tex}
\input{_tex/rules-dataize.tex}
\end{mdframed}
\capt{Dataization rules, generated by phino \iexec{phino --version}\unskip.}
\label{fig:dataization}
\end{figure*}

\begin{theorem}\label{thm:dataization-confluent}
The dataization function is confluent, meaning that
  every normal form can be dataized to at most one datum,
  provided the functions attached to atoms are pure---their result
  depends on nothing beyond the formation, the universe, and the state.
\end{theorem}

\Cref{app:dataization-examples} demonstrates how the dataization function works through examples.

\subsection{Congruence}

\begin{definition}[Congruence]
Two expressions \(n_1\) and \(n_2\) are said to be behaviorally equivalent
  or \textbf{congruent} (denoted by \(\cong\))
  if for any state \(s_1\) and any \(e\):
\begin{phiquation*}
\phinoDataize{n_1}{e}{s_1}{n}{s_2} \quad\text{and}\quad \phinoDataize{n_2}{e}{s_1}{n}{s_2}.
\end{phiquation*}
\end{definition}

Two congruent expressions may be non-equivalent, for example:
\begin{phiquation*}
e_1 = [[ foo -> ?, D> 01-02 ]] \quad e_2 = [[ bar -> ?, D> 01-02 ]]
e_1 \cong e_2 \;\not\to\; e_1 \equiv e_2 {.}
\end{phiquation*}
